\section{Introducción}

En el campo de la cardiologia debido a la complejidad y gravedad de las enfermades se hace de especial importancia el uso de herramientas capaces de leer y analizar datos clinicos de los pacientes de manera que se pueda advertir de posibles riesgos para su salud, sabiendo como esta tarea puede encontrar solucion en las estructuras de datos como grupo vemos de especial interes hacer uso de lo visto en el curso
para clasificar los datos clinicos de un paciente de forma que se puedan leer y analizar de manera eficiente por parte del personal de salud.

Para este fin nuestro proyecto abordara e implementara un algoritmo para la lectura, analisis y graficacion de señales ECG (obtenidas a base de un electrocardiograma) para ser capaces de entregar un reporte con posibles motivos de comportamientos anomalos por medio de algoritmos capaces de predecir patrones.

\subsection{Definicion del proyecto}
La idea central se basa en apoyarse de los temas que se han ido aprendiendo a lo largo del semestre sobre diferentes estructuras de datos para implementarlas en el manejo de los datos del ECG y detectar tanto propiedades como anomalias de este, esta metodologia permitira innovacion y replanteamiento de ideas a largo del proceso de desarrollo al a vez que la adicion de mejoras en cuanto a eficiencia y funcionamiento que facilmente pueden entrar tanto en el contexto del curso como del propio proyecto. 


\subsection{Objetivos}

\begin{itemize}
    \item \textbf{Ser capaz de leer registros de la base de datos MIT-BIH Arrhythmia Database (2 canales), PTBDB (15 canales) y PTB-XL(12 canales) incluyendo los archivos de registro y las anotaciones del cardiologo.} \cite{MIT-BIH} \cite{PhysioNet}
    \item \textbf{Hacer un uso activo de estructuras de datos vistas en el curso capaces de permitir una recepcion y envio eficiente de datos para el procesamiento y filtrado de informacion}
    \item \textbf{Implementar un algoritmo capaz de predecir casos a partir de modelos para entrenamiento de algoritmos, esto se hara por medio del uso de los random forest (los cuales son a su vez una implementacion de arboles de decision binarios)}
    \item \textbf{Entregar un resultado facil de leer para el cardiologo de manera que permita analizar los datos facilmente} 
\end{itemize}